\section*{Impact Statement}

This paper studies how learned world models represent and extrapolate a physical law in a
synthetic environment, and proposes measurement tools (matched-state counterfactuals, nested
inversion, a Bayes-optimal reference) for attributing out-of-distribution failures to a stage of
the model. The direct societal consequences are those of better evaluation: world models are
increasingly used to plan actions for robots and to generate video, and methods that reveal
\emph{why} a model fails outside its training distribution can make such systems safer to deploy
and harder to oversell. The experiments use no human data and no real-world footage. We see no
specific negative consequences beyond those common to interpretability research in general.
